\section{Passing from layer points to the original family class}

% BEGIN X86B1
\paragraph{X86B1. A family bound from eventual member hypotheses.}
For a datum $D$ and its specified norm choice $a$, write
$P(D,a)$ for the conjunction of the eighteen Galois preconditions
and the four ties. Work under the published hypotheses of A85D3.
Fix a family $F$ of positive scale exponent with $\mathrm{TwGood}(F)$.
Suppose there is a real $b$ such that, for every $x\ge b$, every
$h:F.x_1<x$ and every $i\in F.I_x$,
\[
               P\bigl(F.D(x,i),F.\mathrm{ex}(x,h,i).\mathrm{nch}\bigr).
\]
Then $\mathrm{LinkN}(F)$ holds.

Put $s=\max(F.x_1,b)$ and form a family $F'$ with starting
threshold $s$. It has the original index sets at \emph{every} $x$,
including $0$, the original unbounded set of windows, the same $D,L$
and the same numerical constants. To specify every field involving
the new guard $s<x$, supply to the old field its composite
$F.x_1\le s<x$. In particular use these restrictions for the parameter
choice and its comparisons, $\mathrm{parT}$, the size and covering
properties, and $\mathrm{ex}$. All unguarded fields are the given
ones and $1<F.x_1\le s$. Equal proofs of a real inequality give
identical values of every such dependent field. This constructs
the entire $\mathrm{FamP}$ record, including its family coefficient
lists. The same witnesses and cutoff in $\mathrm{TwGood}(F)$ prove
$\mathrm{TwGood}(F')$: its norm choice and its theta datum are
the restrictions of precisely the ones used in the original test.
At every $s<x$, the displayed hypothesis now supplies the member
input of A85D3. Apply it to $F'$, with its type-prime constant $40$.

Let $C,x_0'$ be its output. Take $x_0=\max(x_0',s+1)$ on $F$.
For all $x\ge x_0$ and all $i$, both statements use the same
$D,L$, lifts, $dZ_0$, $\eta$-sum and cyclotomic series.
The $\mathrm{ShaFinOdd}$ and summability implications therefore
give exactly the desired estimate with $C,x_0$. No class membership
of $F'$ is needed in this application of the family theorem.

Consequently, for an original fixed class $\mathcal C$, it suffices
to have this eventual member assertion as a $\mathrm{Closed}$
statement on $\mathrm{Des}(\mathcal C,3)$. Retain its
$\varepsilon_1$ and its transport parameter $K(\varepsilon)$
before taking the family or $C_H$. The definition of the design
supplies $\mathrm{TwGood}$ on the given family. Choose $b$ there,
then apply the preceding argument. Separate closed inputs can be
combined with the minimum of their positive exponent cutoffs,
the pointwise supremum of the four components of their transport
parameters, and the maximum of their member thresholds on this
family. These choices retain the order of all the original quantifiers.
% END X86B1

% BEGIN X86B2
\paragraph{X86B2. Which selected row supplies the norm choice.}
Consider the predicate $P$ of X86B1. Suppose that on the original
selection $\sigma$ of a family class $\mathcal A$ we have the
following row assertion: for every $F\in\mathcal A$, every
$h:F.x_1<x$ and every $i$, the actual object
$O=\mathrm{objAt}(F,x,h,i)$ has $P(F.D(x,i),O.\mathrm{ex}.\mathrm{nch})$.
Then
\[
       \mathrm{Closed}\bigl(D_3(3,\mathcal A),\mathrm{LinkN}\bigr),
       \qquad D_3(3,\mathcal A)=
                  \mathrm{Des}(\mathrm{WithChN}(\mathcal A),3).
\]
Indeed membership projects to $F\in\mathcal A$ and to its twist
condition. The row uses the same $\mathrm{ex}$ and gives X86B1 with
$b=F.x_1$. Thus one may take $\varepsilon_1=1$ and the transport
parameter with all four components zero, uniformly before $F,C_H$.
The family output constants still have their permitted dependence.
In particular this deduction uses no separate bound for primes
$2,3$ on the selected class. The source's combination with a closed
$2,3$ row takes its outer $\varepsilon_1,K$ and applies the
family theorem using the member row alone. A85D3 already includes
the necessary arithmetic estimates on the qualifying family.

Here is a sufficient way to verify the row premise while keeping
the existing selection fixed. For each of its members require a
witness $q$ over the layer and the equality in the dependent sum
\[
        (F.D(x,i),O)=(q.w.D(q.Z),q.O).
\]
Require that its base and frame are within the range of the
layer-point member assertion, with the specified norm core, default
extension, row extension and column pin there. Evaluate the assertion
at this witness and substitute the single displayed equality:
it transports its norm argument to $O.\mathrm{ex}.\mathrm{nch}$.
By X86B1 the witnesses may also be provided for every $i$ after a
common threshold depending on the family.

The distinction between this hypothesis and the bare assertion
at layer points is visible in the source. The latter is an eventual
statement in the base context. It produces a permitted additional
condition on the slot, called \texttt{membershipLocation}, whose
content is its universally quantified conclusion at that base.
On the conjunction of this condition with the good built slot the
construction takes a built witness and its goodness, then the
division lattice and height of the selected point. Transport of the
witness along that division change produces $q$ with the displayed
equality. Its norm extension projects to the default extension,
and the witness itself gives the row extension and column pin.
This is the pointwise substitution used by the source's linked
row theorem. The linked conclusion quantifies existentially over
such an extra slot condition; the row in turn specifies an eventual
predicate $E$ and applies to selections realizing $E$. For a
previously fixed $\mathrm{Des}(\mathcal C,3)$, the specialization
must therefore supply the matching witnesses on its original
families, or show that its existing slot and selector realize
these predicates there. Inclusion of the original class in the
verified domain suffices; changing the original selected class to
a smaller class does not establish the estimate for its other
members. This describes the remaining family binding when the
remaining upper multiplicity comparison of F86D1--2 has been supplied.
% END X86B2
